\input{template/config}
\renewcommand{\thetable}{\arabic{table}}
\renewcommand{\tablename}{Tabla}
\captionsetup[table]{
	format=plain,
	labelfont=bf,
	labelsep=colon
}


La metodología del proyecto se organizó en siete fases consecutivas que abarcan desde la recolección y preparación de los datos hasta la implementación del sistema de información. El enfoque adoptado se basa en el concepto de valor agregado institucional, donde la red neuronal cumple la función de estimar el desempeño esperado de los estudiantes en la prueba Saber Pro. A partir de esta estimación, el valor agregado se calcula como la diferencia entre el resultado realmente obtenido por el estudiante y el resultado que el modelo estimaba para él.

\subsection{Recolección e integración de datos}
Para el estudio se utilizaron los registros de las pruebas Saber 11 (2010-2024) y Saber Pro (2014-2023), almacenados en una base de datos PostgreSQL compuesta por 25 tablas de resultados, correspondientes a quince años de Saber 11 y diez años de Saber Pro. La integración de la información se realizó directamente en el motor de base de datos, donde se consolidaron las tablas de cada prueba y se relacionaron los registros de los estudiantes que presentaron Saber Pro con sus respectivos resultados en Saber 11 mediante una tabla de llaves. Como resultado de este proceso, se obtuvieron aproximadamente $1,32$ millones de pares de estudiantes emparejados, a partir de cerca de $6,05$ millones de registros de Saber 11 y $1,88$ millones de registros de Saber Pro. Los registros que lograron emparejarse correctamente conformaron la muestra utilizada en el estudio. Posteriormente, el procesamiento y análisis de los datos se llevó a cabo en Python utilizando bibliotecas como PySpark, Pandas, NumPy, Scikit-learn y TensorFlow/Keras. Debido al volumen y a la forma en que se gestionó la información, fue posible ejecutar todo el proceso en memoria, sin necesidad de recurrir a infraestructuras de cómputo distribuido.

\subsection{Depuración y tratamiento de registros atípicos}
Durante el análisis se identificaron registros de puntajes considerados inusuales, caracterizados por presentar un valor de cero en los módulos. Estos registros se conservaron en la base de datos con el fin de preservar la integridad de la información original; sin embargo, no se tuvieron en cuenta durante los análisis estadísticos tampoco durante el entrenamiento de los modelos, debido a que un valor de cero en un puntaje no representa un desempeño real y su inclusión distorsionaría las medidas descriptivas, las correlaciones y la actualización de los parámetros del modelo. El origen de estos registros se asocia a las pruebas Saber Pro anteriores al año 2016, periodo en el que se utilizaba un método de calificación distinto al vigente; por ello, dichos registros fueron excluidos del análisis, no de la base de datos, para no mezclar formatos de calificación incompatibles, decisión que constituye un supuesto metodológico adoptado durante la depuración. De manera análoga, se descartaron de los análisis las filas con datos ausentes en los módulos análogos que actúan como variables de entrada o de salida. En sentido contrario, se mantuvo una categoría "-A1" del módulo de Inglés: si bien el Marco Común Europeo de Referencia (MCER) clasifica la competencia en seis niveles (A1 a C2), el conjunto de datos incluye un nivel adicional "-A1" que corresponde a un dominio prácticamente nulo del idioma, situación frecuente en la población colombiana; por su elevada frecuencia, excluir esta categoría habría sesgado el análisis del módulo, razón por la cual se conservó.

\subsection{Correspondencia de competencias y cálculo del puntaje global}
Dado que las competencias evaluadas en Saber 11 y Saber Pro no tienen exactamente la misma denominación, fue necesario definir una correspondencia entre aquellos módulos que evalúan habilidades equivalentes en ambas pruebas (véase Tabla~\ref{tab:correspondencia}). A partir de este ejercicio se identificaron cuatro pares comparables: la competencia lógico-cuantitativa, la lectoracomunicativa, la ciudadana-cívica y la de Inglés. El módulo de Ciencias Naturales de Saber 11 no se incluyó dentro del análisis comparativo, ya que no cuenta con un componente equivalente en la prueba Saber Pro. Respecto al puntaje global, la mayoría de los registros lo traían ya calculado por el ICFES; la excepción son los años 2010 a 2013 de Saber 11, para los cuales no existía un puntaje global dado que en ese periodo el ICFES no lo calculaba, situación que responde a un cambio histórico en la estructura del examen y no a una ausencia de datos. El ICFES obtiene el global de Saber 11 como un promedio ponderado de la siguiente forma (Ecuación~\eqref{eq:global_saber11_icfes}):
\begin{equation}
\frac{(3 \cdot MA) + (3 \cdot LC) + (3 \cdot SC) + (3 \cdot CN) + (1 \cdot IN)}{13} \cdot 5
\label{eq:global_saber11_icfes}
\end{equation}
donde $MA =$ puntaje matemáticas, $LC =$ puntaje lectura crítica, $SC =$ puntaje ciencias sociales/ciudadanas, $CN =$ puntaje ciencias naturales y $IN =$ puntaje inglés; mientras que el de Saber Pro lo calcula como un promedio simple tal que (Ecuación~\eqref{eq:global_saberpro_icfes}):
\begin{equation}
\frac{RC + CE + LC + CC + IN}{5}
\label{eq:global_saberpro_icfes}
\end{equation}
donde $RC =$ puntaje razonamiento cuantitativo, $CE =$ puntaje comunicación escrita, $LC =$ puntaje lectura crítica, $CC =$ puntaje competencias ciudadanas y $IN =$ puntaje inglés \cite{ICFES_FichaMetodologica2022} \cite{MEN_SINEB_Aprendizaje2023}. El menor peso del Inglés en Saber 11 responde a un criterio de equidad: el nivel socioeconómico condiciona la calidad de la educación recibida y este efecto es especialmente marcado en Inglés, donde los estudiantes de instituciones privadas o de mayores recursos obtienen puntajes notablemente superiores. Sobre esta base, en el presente trabajo se calculó un puntaje global propio para ambas pruebas, replicando la lógica del ICFES pero considerando únicamente los cuatro módulos con análogo, reestructurando el modelo matemático de la forma tal que no se generen discrepancias significativas que afecten las métricas. (Ecuaciones~\eqref{eq:global_saber11_propio} y \eqref{eq:global_saberpro_propio}):
\begin{equation}
\frac{(3 \cdot MA) + (3 \cdot LC) + (3 \cdot SC) + (1 \cdot IN)}{10} \cdot 5
\label{eq:global_saber11_propio}
\end{equation}
\begin{equation}
\frac{RC + LC + CC + IN}{4}
\label{eq:global_saberpro_propio}
\end{equation}
Bajo este argumento, tanto para Saber 11 como para  Saber Pro, se realiza la siguiente propuesta.

El puntaje global de Saber Pro se calcula también como promedio ponderado con menor peso para el Inglés y no como promedio simple, dado que la asociación entre nivel socioeconómico y dominio del idioma se mantiene en la educación superior. Para sintetizar el proceso ejecutado, la Tabla~\ref{tab:correspondencia} expone las relaciones identificadas y su respectiva segmento asociado como competencia general.

\begin{table}[!ht]
	\centering
	\begin{adjustbox}{max width=\textwidth}
		\footnotesize
		\begin{tabularx}{\textwidth}{|>{\raggedright\arraybackslash}X|
				>{\raggedright\arraybackslash}X|
				>{\raggedright\arraybackslash}X|}
			\hline
			\textbf{Saber 11} &
			\textbf{Saber Pro} &
			\textbf{Competencia general} \\ \hline
			
			punt\_matematicas\_norm &
			mod\_razona\_cuantitat\_punt\_norm &
			Competencia lógico-cuantitativa \\ \hline
			
			punt\_lenguaje\_norm (2010--2013), punt\_lectura\_critica\_norm (2014 en adelante) &
			mod\_lectura\_critica\_punt\_norm &
			Competencia lectora-comunicativa \\ \hline
			
			punt\_ciencias\_sociales\_norm (2010--2013), punt\_sociales\_ciudadanas\_norm (2014 en adelante) &
			mod\_competen\_ciudada\_punt\_norm &
			Competencia ciudadana-cívica \\ \hline
			
			punt\_ingles\_norm &
			mod\_ingles\_punt\_norm &
			Competencia en Inglés \\ \hline
			
			punt\_global\_norm &
			punt\_global\_norm &
			Puntaje Global \\ \hline
			
			punt\_global\_calc\_norm &
			punt\_global\_calc\_norm &
			Puntaje Global Calculado \\ \hline
			
		\end{tabularx}
	\end{adjustbox}
	\caption{Tabla de correspondencia de pares análogos entre Saber 11 y Saber Pro.}
	\label{tab:correspondencia}
\end{table}

\subsection{Normalización y escalado}
El ICFES no entrega los puntajes normalizados, sino en las escalas propias de cada prueba, que no son directamente comparables entre Saber 11 y Saber Pro. Por ello, durante el desarrollo del sistema se aplicó una normalización Min-Max al rango $[0, 1]$, con el fin de llevar ambos exámenes a una escala común que habilitara la comparación de competencias (criterio de dimensionalidad). Es preciso advertir que el ICFES reformó la metodología de la prueba Saber 11 en 2014; en consecuencia, los años anteriores no son comparables con los posteriores, tanto porque antes de esa fecha no existía un puntaje global como porque el método de calificación de los módulos específicos era distinto del actual. Esto significa que, aunque la escala de puntajes sea similar, los resultados obtenidos entre 2010 y 2013 no pueden compararse directamente con los registrados a partir de 2014. Si bien el ICFES aplica estas evaluaciones desde años anteriores, para este estudio solo se dispuso de información de Saber 11 desde 2010. A pesar de esta diferencia, los puntajes correspondientes al periodo 2010-2013 fueron normalizados utilizando el mismo procedimiento aplicado al resto de los datos. Esta decisión se adoptó y se reconoce de manera explícita como un supuesto metodológico del estudio. Por otra parte, las variables utilizadas como entrada en los modelos predictivos fueron estandarizadas mediante el método de puntaje z antes del proceso de entrenamiento. Asimismo, los parámetros de escala y los valores empleados para imputar los datos faltantes de las variables socioeconómicas, utilizando la mediana como criterio de reemplazo, se calcularon únicamente a partir del conjunto de entrenamiento. De este modo, se evitó que información procedente de los conjuntos de validación y prueba influyera en el ajuste.

\subsection{Partición temporal}
De cara al entrenamiento de la RNA, La partición de los datos no se realizó de forma aleatoria sino por cohortes según el año de presentación de Saber 11: entrenamiento con las cohortes 2010–2016, validación con 2017 y prueba con 2018. Se excluyeron las cohortes de Saber 11 posteriores a 2018 por no contar aún con su correspondiente Saber Pro dado que estos puntajes están en proceso. Sobre el conjunto se calcularon medidas descriptivas y se analizó la correlación entre las competencias iniciales de Saber 11 y las finales de Saber Pro.

\subsection{Sistema de formas y variables de entrada}
La estimación se abordó bajo dos configuraciones de entrada, denominadas formas, con el fin de cuantificar el aporte del contexto socioeconómico. La Forma 1, utilizada como modelo de referencia, considera únicamente los cuatro módulos equivalentes de la prueba Saber 11: razonamiento cuantitativo, lectura crítica, competencias ciudadanas e inglés. Además, en los modelos asociados a razonamiento cuantitativo y competencias ciudadanas se incluyeron tres variables relacionadas con el tiempo: el año en que el estudiante presentó Saber 11, el año de presentación de Saber Pro y la diferencia, en años, entre ambas evaluaciones. Por su parte, la Forma 2 amplía el conjunto de variables de entrada hasta un total de veinte. Además de los cuatro módulos académicos, incorpora ocho variables de carácter socioeconómico, siete indicadores relacionados con la ausencia de información y una variable que identifica la institución objetivo. A diferencia de la Forma 1, esta configuración no incluye las variables temporales.
Las ocho variables socioeconómicas provienen del registro de Saber 11 y describen el contexto del estudiante: el carácter bilingüe del colegio, su área de ubicación (urbana o rural), su naturaleza (oficial o privada), la jornada escolar, el género del estudiante, el estrato de la vivienda y el nivel educativo de la madre y del padre. Dado que algunas presentaban valores ausentes, se imputaron por la mediana calculada solo con el conjunto de entrenamiento y, por cada variable con ausencias apreciables, se construyó un indicador de ausencia: una variable booleana que toma el valor uno cuando el dato original faltaba y cero cuando estaba presente, de modo que el modelo pueda distinguir un valor imputado de uno real (siete en total, pues el género, con menos del cinco por ciento de ausencias, no requirió indicador). El indicador de institución objetivo es, a su vez, una variable booleana que señala si el estudiante presentó Saber Pro en la Universidad de San Buenaventura, lo que permite filtrar y comparar posteriormente el desempeño de la institución frente al total. El contraste entre ambas formas constituye en sí mismo un resultado, pues la diferencia de error entre la Forma 1 y la Forma 2 cuantifica el aporte del contexto socioeconómico a la estimación.

\subsection{Arquitectura y entrenamiento del modelo de valor agregado}
El motor predictivo del sistema es una red neuronal artificial de tipo perceptrón multicapa. En lugar de una sola red que estime los cuatro módulos a la vez, se entrenó una red independiente por cada competencia y por cada forma, de modo que cada una optimiza un único objetivo sin que el error de un módulo interfiera con el de los demás. 
Con el propósito de reducir el efecto de la variabilidad asociada a la inicialización aleatoria de las redes neuronales, cada modelo se entrenó tres veces utilizando semillas diferentes $(42, 123$ y $777)$. La predicción final para cada estudiante se calculó como el promedio de los resultados obtenidos en estas tres ejecuciones. Esta estrategia de ensamble permite obtener estimaciones más estables al disminuir la variabilidad de las predicciones sin modificar el sesgo del modelo. Por esta razón, las métricas reportadas corresponden al desempeño del ensamble completo evaluado sobre el conjunto de prueba y no al promedio de las métricas obtenidas en cada ejecución individual, ya que este último enfoque no refleja adecuadamente los beneficios del promediado. 
Todas las redes utilizaron la misma arquitectura base, compuesta por tres capas ocultas de tamaño decreciente. La primera cuenta con $64$ neuronas e incorpora normalización por lotes y una función de activación LeakyReLU. La segunda está formada por $32$ neuronas con activación ReLU y una tasa de abandono (dropout) del $20\%$. La tercera contiene $16$ neuronas con activación ReLU, mientras que la capa de salida está integrada por una única neurona lineal encargada de generar el puntaje continuo de cada módulo. Los pesos iniciales se definieron mediante el método de He (Kaiming), y el entrenamiento se realizó utilizando el optimizador Adam con una tasa inicial de aprendizaje de $0,001$ y la función de pérdida de error cuadrático medio. El proceso se ejecutó con lotes de $512$ observaciones y un máximo de cien épocas de entrenamiento. Para reducir el riesgo de sobreajuste se implementaron mecanismos de detención temprana, que restauran automáticamente los pesos correspondientes al mejor resultado observado en el conjunto de validación después de quince épocas sin mejora. Adicionalmente, la tasa de aprendizaje se ajustó de manera dinámica cuando el modelo dejaba de mostrar avances significativos durante el entrenamiento. Con el fin de garantizar la reproducibilidad de los experimentos, el estado del generador aleatorio se reinició antes de construir cada red. De esta manera, fue posible comparar los resultados obtenidos con distintas semillas bajo condiciones equivalentes. 
Para cada forma se calculó una única estandarización a partir del conjunto de entrenamiento y posteriormente se aplicó a las cuatro redes correspondientes. Cuando se utilizaron variables temporales, estas se incorporaron al final del vector de entrada. Este procedimiento permitió que los módulos que no empleaban dichas variables mantuvieran exactamente la misma estructura de datos base, preservando la coherencia del puntaje global reconstruido a partir de los cuatro módulos evaluados. La configuración denominada 'V3\_5' surgió a partir de un proceso de experimentación incremental. Sobre una arquitectura inicial se evaluaron distintos cambios de forma independiente, incluyendo nuevas variables, funciones de pérdida robustas, variables temporales y modelos especializados para cada módulo. Únicamente se conservaron las modificaciones que aportaron mejoras verificables en el desempeño, mientras que aquellas que no generaron beneficios o produjeron resultados inferiores fueron descartadas. Como resultado, la única especialización que se mantuvo fue la incorporación de tres variables temporales (año de presentación de Saber 11, año de presentación de Saber Pro y número de años transcurridos entre ambas pruebas) en los módulos de razonamiento cuantitativo y competencias ciudadanas de la Forma 1. Los demás módulos, tanto en la Forma 1 como en la Forma 2, conservaron la configuración base utilizando el error cuadrático medio como función de pérdida.
Como modelos de referencia, frente a los cuales se contrasta la red, se construyeron una regresión lineal y un esquema de segmentación K-means bajo la estrategia agrupar-y-predecir descrita en el Marco Conceptual. La técnica de árboles de decisión, contemplada inicialmente, no se implementó y se plantea como trabajo futuro.

\subsection{Evaluación y despliegue del sistema}
El desempeño de los modelos se evaluó utilizando el error absoluto medio (\textit{MAE}), la raíz del error cuadrático medio (\textit{RMSE}) y el coeficiente de determinación ($R^2$) para cada competencia analizada. Para calcular el puntaje global predicho, primero se desnormalizaron los resultados de cada módulo, luego se aplicó la fórmula de ponderación establecida y, finalmente, el resultado obtenido se normalizó nuevamente. A partir de estos valores, el valor agregado institucional se estimó para cada competencia y para cada año como la diferencia entre el desempeño observado y el desempeño esperado. Los resultados obtenidos fueron integrados en un sistema de información desarrollado con Dash sobre una base de datos PostgreSQL. La herramienta incorpora visualizaciones con líneas de tendencia que permiten comparar los resultados de Saber 11 y Saber Pro para un mismo grupo de estudiantes, así como una matriz diseñada para medir el aporte de la institución en el desarrollo de las competencias evaluadas. Para su puesta en funcionamiento, el sistema fue preparado para ejecutarse en un servidor Ubuntu utilizando Gunicorn, Nginx y systemd. La población objeto de estudio estuvo conformada por los estudiantes de la Universidad de San Buenaventura que cuentan con registros tanto en Saber 11 como en Saber Pro, mientras que la muestra correspondió a los pares de estudiantes que pudieron ser emparejados correctamente durante el proceso de integración de datos.